\documentclass[11pt]{article}
\usepackage[a4paper,margin=1in]{geometry}
\usepackage{amsmath,amssymb,amsthm}
\usepackage{booktabs}
\usepackage{graphicx}
\usepackage{xcolor}
\usepackage{hyperref}
\hypersetup{colorlinks=true,linkcolor=blue,urlcolor=blue,citecolor=blue}
\newcommand{\rt}[1]{\texttt{#1}}
\newcommand{\fn}[1]{\texttt{#1}}
\newcommand{\cmark}{\ensuremath{\checkmark}}
\newcommand{\xmark}{\ensuremath{\times}}
\providecommand{\ket}[1]{|#1\rangle}
\newcommand{\wcc}{\mathrm{wcc}}
\newcommand{\att}{\mathrm{att}}
\newcommand{\scc}{\mathrm{scc}}
\newcommand{\ifgraphic}[2]{\IfFileExists{#1}{\includegraphics[width=#2]{#1}}%
  {\fbox{\parbox{#2}{\centering\small figure \texttt{\detokenize{#1}} pending
  (regenerate after the sweep)}}}}
\newcommand{\iftable}[1]{\IfFileExists{#1}{\input{#1}}%
  {\fbox{\parbox{0.9\textwidth}{\centering\small table
  \texttt{\detokenize{#1}} pending (regenerate after the sweep)}}}}

\theoremstyle{plain}
\newtheorem{theorem}{Theorem}
\newtheorem{proposition}{Proposition}
\theoremstyle{definition}
\newtheorem{definition}{Definition}

\title{\textbf{Report R9 --- Sectors and basins: the weak-component
decomposition over the rule space}\\
\large Tier 1e, open boundaries}
\author{QCA Sector Analysis project (Tier 1)}
\date{2026-07-29 \quad\small(\fn{tier1e\_version} \texttt{1e.1}, \rt{obc0})}

\begin{document}
\maketitle

\begin{abstract}
The dissipative pipeline has so far reported terminal strongly connected
components --- ``recurrent classes'' --- and their basins. Those are the right
object for the \emph{measured} circuit and the wrong object for a sector axis:
they do not partition the basis, they are monitored-relative, and the Tier-2
wall grammar predicts sectors rather than attractors. This report adds the
weakly connected component (WCC) as a first-class observable, which fixes all
three at once: $\operatorname{span}(\text{WCC})$ is an enclosure of the channel,
invariant under every Kraus operator and therefore exact for the monitored and
the unmonitored dynamics alike; WCCs partition the basis; and they are precisely
the minimal projectors of the \emph{diagonal} part of the commutant. We sweep
all $256$ rules at open boundaries, verify the sum rule and the cross-tier
consistency assertions, and test the exclusion curve $ab\ge2$ that a partition
forces on the sector map. Two negative results are recorded prominently: the
Tier-2 prediction $a_\wcc\ge\varphi$ for rule $28$ is \emph{false} --- the sector
count obeys an exact recurrence with base the plastic number
$\rho=1.32472$ --- and the basin sum rule turned out to be unanswerable from the
Tier-1a archive for $76$ units, which we fixed by recomputing them
independently. This report covers \rt{obc0} only; \rt{pbc} follows separately.
\end{abstract}

\tableofcontents

\section{Three tiers, and why the WCC is the sector-level object}
\label{sec:tiers}

\subsection{Vocabulary}

The single most important thing this report changes is naming. Throughout the
project from here on:

\begin{center}
\begin{tabular}{ll}
\toprule
\textbf{sector} & weakly connected component; exact for \emph{both} experiments\\
\textbf{monitored attractor} & terminal SCC; a property of the measured circuit\\
\textbf{channel attractor} & block of the fixed-point algebra (Tier 1b / 1d)\\
\bottomrule
\end{tabular}
\end{center}

\noindent A terminal SCC is never a ``sector''; a WCC is never an ``attractor''.
Every table and figure caption below states which experiment its numbers
describe.

\subsection{Why terminal SCCs cannot carry a sector axis}

\begin{description}
\item[(P1) They do not partition.] Transient states belong to no recurrent
  class, so $\sum_k|A_k|\ll 2^N$ and no sum rule constrains the multiset. The
  unitary sectors \emph{do} partition, so the existing figures put two
  different kinds of object on one axis. Rule $22$ at $N=12$ is the extreme
  case: three monitored attractors, each a single state, covering $3$ of
  $4096$ states, against two sectors covering all $4096$.
\item[(P2) They are monitored-relative.] The transient/recurrent split is a
  property of the measured chain; the unmonitored channel can keep immortal
  weight on transient states (R5 rule 22, R7).
\item[(P3) They are not what the analytics predicts.] The Tier-2 wall grammar
  and transfer matrices give sector counts.
\end{description}

\subsection{What a WCC is}

\begin{definition}
The one-cycle support graph has nodes $x\in\{0,1\}^N$ and an edge $x\to y$
whenever $\langle y|\Phi_C(|x\rangle\langle x|)|y\rangle\neq0$. A \emph{sector}
is a connected component of this graph read as undirected.
\end{definition}

\begin{proposition}[enclosure]
$\operatorname{span}(\text{sector})$ is invariant under every Kraus operator of
$\Phi_C$, hence invariant for the monitored and the unmonitored dynamics alike.
\end{proposition}

\noindent The commutant reading is the one that connects this tier to R8. A
diagonal operator $D$ commutes with the dynamics iff $D_x=D_y$ whenever the
support graph has an edge $x\to y$, i.e.\ iff $D$ is constant on components. So
\begin{equation}
\dim\{\,D \text{ diagonal}:[D,\Phi_C]=0\,\}\;=\;\#\text{components}\;=\;n_\wcc,
\label{eq:diagcommutant}
\end{equation}
and the graph analysis computes \emph{exactly} the diagonal part of the
commutant --- no more. That is why it cannot see reflection, or any other
symmetry whose eigenvectors are superpositions rather than basis states
(R8~\S10.3: at $N=11$ the diagonal part is $7$-dimensional inside a commutant of
dimension $23\,360$).

\subsection{The algorithm, and one deliberate deviation}

A single streaming pass over $\operatorname{succ}(x)$ for
$x=0\ldots2^N-1$, union-find with union by size and path compression, every edge
symmetrised by construction ($\mathrm{union}(x,y)$ for each $y\in
\operatorname{succ}(x)$). No adjacency storage, no DFS stack. This is the code
path the unitary sweep already used, now factored into
\fn{graph/wcc.py::weak\_components} and shared by both callers.

\paragraph{Deviation, and why.} The Tarjan sweep aborts a unit as soon as one
component exceeds $f_{\rm erg}2^N$. The WCC pass does \emph{not}: aborting
leaves an incomplete partition, which would break the sum rule A1 and destroy
the exact anchors the exclusion-curve test needs --- rule $51$ has to come out
as one sector of size exactly $2^N$, not as ``something exceeded half the
space''. Union-find is cheap enough to always finish, so \rt{ergodic} here is a
\emph{classification} derived from $D_{\max}$, and the exponential saving lives
in the sweep instead. A second guard: a rule is not written off as ergodic on
the strength of a tiny chain, since at $N=6$ almost anything connects past half
of $64$ states; the verdict must hold at two consecutive $N$ and never below
$N=10$.

\section{The sweep}
\label{sec:sweep}

All $256$ rules at \rt{obc0}, swept with $N$ in the \emph{outer} loop so that a
run stopped against a wall-clock budget still leaves uniform coverage --- every
rule with the same $N$ range and therefore comparable fits. Records live in
their own store \fn{results\_tier1e/}, keyed on $(\text{rule},N,\rt{bc})$ and
versioned \texttt{1e.1}.

\paragraph{A note on versioning.} The task asks for the shared
\fn{ENGINE\_VERSION} to be bumped to \texttt{1e.1}. Doing that would make
\fn{has\_unit()} false for all ${\sim}6000$ archived Tier-1a records, so the next
sweep would recompute the lot --- which the same section explicitly forbids. The
version bump therefore lives on the new tier, exactly as Tier 1d does. Nothing
in Tier 1a/1b/1d is recomputed or invalidated.

\begin{center}
\iftable{tab_r9_coverage_obc0.tex}
\end{center}

\noindent Coverage is uniform: all $256$ rules at every $N=6\ldots16$, so the
descriptors of \S\ref{sec:hyper} are all fitted on the same domain. A targeted
set of $30$ diagnostic rules was carried to $N=17$--$18$ afterwards; the headline
map deliberately does \emph{not} use those extra points, since a rule fitted on a
longer window is not comparable with one fitted on the uniform window.

\section{Ground truths}
\label{sec:truth}

All of the following are regression-tested in \fn{tests/test\_wcc.py}.

\begin{itemize}
\item \textbf{Rule 204} (\rt{IIII}): $2^N$ sectors of size $1$; $a=2$, $b=1$.
\item \textbf{Rule 51} (\rt{VVVV}): one sector of size $2^N$; $a=1$, $b=2$.
\item \textbf{Rule 150} \rt{obc0}: sectors are the domain-wall-number
  eigenspaces, $|S_w|=\binom{N+1}{w}$ on even $w$. The computed values reproduce
  the reference series for $N=8\ldots17$ exactly --- $\#$sectors
  $5,6,6,7,7,8,8,9,9,10$ and $D_{\max}$ $126$, $210$, $462$, $924$, $1716$,
  $3003$, $6435$, $12870$, $24310$, $43758$.
\item Two of the task's regression targets are \rt{pbc}-specified (rule $156$'s
  Lucas law and rule $22$'s attractor count). Both are implemented and pass in
  the test suite, but they are \rt{pbc} statements and are therefore discussed in
  the \rt{pbc} report, not here.
\end{itemize}

\paragraph{A naming trap worth recording.} Wolfram numbers are not classical ECA
rules in this encoding. $W_{90}$ has symbol tuple \rt{DEED} and is
\emph{dissipative}; assertion A2 must not be applied to it. Its $N=8$ \rt{obc0}
numbers make the P1 point cleanly: sector sizes $112,112,16,16$ summing to
$256$, against terminal-SCC sizes $7,7,1,1$ summing to $16$.

\section{Rule 28: the wall-transparency prediction fails}
\label{sec:r28}

The task predicts $a_\wcc\ge\varphi$ for rule $28$ (\rt{IIVD}), on the grounds
that it inherits the $01$-wall grammar of rule $156$, and asks for a prominent
flag if that is violated. \textbf{It is violated.} The \rt{obc0} sector count is
\[
n_\wcc(N)=11,15,20,27,36,48,64,85,113,150,199 \qquad (N=6\ldots16),
\]
which satisfies the exact integer recurrence
\begin{equation}
a_n=a_{n-1}+a_{n-2}-a_{n-4},
\label{eq:r28}
\end{equation}
verified out of sample on every available term. Its characteristic polynomial
factors as $x^4-x^3-x^2+1=(x-1)(x^3-x-1)$, so the growth base is the
\textbf{plastic number}
\[
\rho=1.324717957\ldots\;<\;\varphi=1.618033988\ldots,
\]
the real root of $x^3=x+1$. This is an exact algebraic base, not a fit, so the
gap to $\varphi$ cannot be attributed to a short series or to a parity artifact.

What this falsifies is the \emph{transparency} step of the argument, not the
wall grammar itself: rule $28$ replaces one of rule $156$'s symbols with a reset
\rt{D}, and the reset evidently removes wall configurations from circulation
rather than merely relabelling them. Quantifying which configurations are lost
is a Tier-2 question and is listed in \S\ref{sec:hooks}.

\section{The hyperbola}
\label{sec:hyper}

\begin{theorem}[exclusion curve]
Sectors partition the basis, so with $K\sim a^N$ sectors and $D_{\max}\sim b^N$,
\begin{equation}
K\,D_{\max}\;\ge\;\sum_k D_k\;=\;2^N
\qquad\Longrightarrow\qquad
ab\;\ge\;2,
\label{eq:hyper}
\end{equation}
with equality iff the sector sizes are all comparable. No rule may lie below
$b=2/a$.
\end{theorem}

\noindent This is a theorem about \emph{partitions}. It therefore constrains the
sector map and says nothing whatever about the monitored-attractor map, where
$a_\att b_\att<2$ is expected and the deficit $2-a_\att b_\att$ measures
transient dominance (V4). We assert nothing there and report the deficit as an
observable.

\subsection{Anchors}

\begin{center}
\iftable{tab_r9_anchors_obc0.tex}
\end{center}

\noindent Rules $204$, $51$ and $150$ sit on the curve to six decimals, as they
must, and rule $156$ sits above it at exactly $\varphi\cdot4^{1/5}=2.1350$. How
the rate behind that number is obtained --- and why it must not be read off a
BIC fit --- is \S\ref{sec:rates}.

\subsection{The check that carries the validation}

The asymptotic statement $ab\ge2$ is a corollary about \emph{bases}, and it only
means anything when both series are exponential: if the sector count is
polynomial its base is $1$ by convention, the product silently drops the
polynomial factor, and the inequality degenerates to $b\ge2$ approached from
below at finite $N$. So the validation is carried instead by the theorem in its
original, fit-free form,
\begin{equation}
n_\wcc(N)\cdot D_{\max}(N)\;\ge\;2^N\qquad\text{for every }N,
\label{eq:finite}
\end{equation}
an immediate corollary of the A1 sum rule involving no fitting anywhere.

\begin{center}
\fbox{\parbox{0.86\textwidth}{\centering
\eqref{eq:finite} holds for \textbf{$256/256$ rules at every $N$}, and it is
\emph{tight}: the smallest value of $n_\wcc D_{\max}/2^N$ anywhere in the sweep
is exactly $1.0000$ (rule $0$ at $N=6$).}}
\end{center}

\subsection{Verdicts in the base plane}

\begin{center}
\iftable{tab_r9_verdicts_obc0.tex}
\end{center}

\noindent The base-plane check is built so that a wide uncertainty band can never
launder a sub-$2$ product into a silent pass. Four things are recorded
separately: the raw product, whether it falls below $2$ at all, whether the
shortfall sits inside the propagated leave-one-out spread, and whether either
series is sub-exponential so that the test is degenerate by construction. Every
rule with $ab<2$ is listed in full, as the task requires.

\begin{center}
\iftable{tab_r9_below_obc0.tex}
\end{center}

\noindent \textbf{There are no violations.} Every remaining sub-$2$ rule
($6$, $14$, $20$, $74$, $84$, $88$, $229$) has a \emph{polynomial} sector count
--- for instance $W_{88}$ has $n_\wcc=3,3,4,4,4,5,5,5,6,6,6$, i.e.\ $\sim N/2$ ---
so $a=1$, and the degenerate form of \eqref{eq:hyper} demands $b\ge2$, which a
finite-$N$ fit necessarily approaches from below ($1.92$--$1.94$ here). Their
$D_{\max}/2^N$ decays slowly enough to be consistent with a power-law prefactor
but not cleanly enough for the test of \S\ref{sec:rates} to certify it; the
theorem, not the fit, gives their base. Two further rules, $W_{90}$ and
$W_{165}$, have genuinely irregular series
($n_\wcc=2,1,4,4,2,6,2,6,12,1,20$) that no growth law describes; they are
classified \emph{irregular} and omitted from the base plane rather than placed at
a fictitious coordinate.

\subsection{Where the rate estimate comes from}
\label{sec:fitting}
\label{sec:rates}

R2~\S3 already documented the trap that this sweep first walked into, and it is
worth restating because it decides four of the numbers above. The BIC model
$\ln y=c+\alpha\ln N+\kappa N$ is the right tool for the growth \emph{class} and
for the sub-leading power, but its $\alpha\ln N$ term absorbs part of the growth,
so $\kappa$ is a \emph{biased} rate estimator. On rule $156$'s $D_{\max}$ it
returns $b=1.2554$ against a true $4^{1/5}=1.31951$ --- far enough off to look
like a different constant, which is exactly how it first appeared here. Rates are
therefore taken, in order of priority:

\begin{enumerate}
\item from an \textbf{exact integer recurrence} where one exists;
\item from an \textbf{analytic derivation} where R2 or R8 supplies one;
\item otherwise from the \textbf{two-parameter fit} $\ln y=c+\kappa N$, which
  cannot trade rate against power.
\end{enumerate}

\noindent The two-parameter fit has the mirror-image bias, and the binomial rules
expose it: with no $\alpha$ to absorb an $N^{-1/2}$ prefactor it pushes the base
\emph{below} $2$ ($1.9244$ for $W_{134}$, whose $D_{\max}$ is exactly rule $150$'s
central binomials). Neither fit is trustworthy alone, so the discriminating
question is asked directly --- is the residual $D_{\max}/2^N$ better explained by
a power of $N$ or by an exponential in $N$? When the power wins, the base is
exactly $2$ and $\alpha$ comes from that fit. That recovers $b=2$ for the whole
binomial family and puts them on the curve beside rule $150$.

The comparison is made on the full series and not only on its tail, and it is
refused when the series falls back down or when the fitted $\alpha$ is too steep
to be a prefactor power ($|\alpha|>8$). All three guards were added on
2026-07-31; the first changed eight of this report's $D_{\max}$ descriptors and
is documented where they are listed (\S\ref{sec:openfrag}), the other two change
nothing here and matter for R10.

\subsection{Two bounds that the fitter needs and the theorem supplies}
\label{sec:inconclusive}

Since $1\le n_\wcc,D_{\max}\le 2^N$, both bases lie in $[1,2]$. A fit outside
that window contradicts a theorem, so the theorem wins: $11$ descriptors were
clamped, each recorded, never silently. The clamp is exactly what the binomial
rules need --- $W_{134}$'s $D_{\max}$ is $924,1716,3003,6435,12870,24310$, i.e.\
precisely rule $150$'s central binomials $\sim2^N/\sqrt N$, and the $M2$ fit
lands at $2.0588$ because the $N^{-1/2}$ prefactor pulls the estimate past the
boundary. Similarly a bounded integer series that has stopped changing is
constant whatever the BIC says: $W_{36}$'s sector count is
$9,9,10,10,\ldots,10$, which the unclamped fitter read as ``exponential with
base $0.95$'' --- below $1$, and impossible for a count.

\paragraph{The hyperbola as a predictor.} Read backwards, \eqref{eq:hyper} is
sharper than any fit, and it called the answer here before the data could. At the
$N\le11$ stage the six rules $140$, $156$, $196$, $198$, $206$, $220$ all had an
\emph{exact} $a=\varphi$ with a $D_{\max}$ series still classified polynomial, so
$ab=\varphi<2$; the theorem already forced $b\ge2/\varphi=1.2361$. With the rate
taken as \S\ref{sec:rates} requires, all six now sit at
\begin{equation}
a\,b\;=\;\varphi\cdot4^{1/5}\;=\;1.61803\times1.31951\;=\;2.1350 ,
\end{equation}
which is exactly the value V3 asks for --- \textbf{above} the curve, with $a$ from
an exact Fibonacci recurrence and $b$ from R2~\S3's room-packing saddle point
rather than from either fit. The six share one $D_{\max}$ series,
$6,9,12,16,20,27,36,48,64,81,108$ over $N=6\ldots16$, so the derivation covers all
of them; the two candidate bases $3^{1/4}=1.31607$ and $4^{1/5}=1.31951$ differ by
$0.26\%$ and no finite series can separate them, which is precisely why the base
is taken from theory.

\begin{figure}[htbp]
\centering
\ifgraphic{../../figures/fig_sector_map_obc0.pdf}{\textwidth}
\caption{\textbf{F1, the two maps on identical axes.} Left: the sector (WCC)
plane, base $a$ of the sector count against base $b$ of the largest sector ---
\emph{exact for both experiments}. Right: the same plane for terminal SCCs ---
\emph{monitored only}. The dotted line is $ab=2$; it excludes the region below it
in the left panel and is only a reference in the right one. Marker area grows
with the number of rules stacked in a cell and the count is printed inside
crowded cells, without which the distribution is invisible: $77\%$ of the sector
map lies in one cell. Colour is by family, with the quantum rules (those carrying
a Hadamard) drawn forward and the V-free rules held back as a baseline. Gold rings and labels mark the eight
V+reset rules that stay exponentially fragmented (\S\ref{sec:openfrag}). The nine
unitary rules land at \emph{identical} coordinates in both panels, which is
assertion A2 made visible.}
\label{fig:f1}
\end{figure}

\section{Sectors against monitored attractors}
\label{sec:contrast}

\subsection{The sector axis is uninformative for the V+reset rules}
\label{sec:pileup}

The left panel of Fig.~\ref{fig:f1} has a blunt message: \textbf{$196$ of the $254$
placeable rules sit at exactly $(a,b)=(1,2)$} --- one sector, of size $2^N$. Broken down by
family:

\begin{center}
\begin{tabular}{lrrr}
\toprule
family & rules & at $(1,2)$ & elsewhere \\
\midrule
unitary (V only) & $16$ & $11$ & $5$ \\
V + reset & $160$ & $\mathbf{150}$ & $10$ \\
V-free baseline & $80$ & $35$ & $43$ \\
\midrule
all & $254$ & $196$ ($77\%$) & $58$ \\
\bottomrule
\end{tabular}
\end{center}

\noindent So $94\%$ of the rules that combine a Hadamard with a reset have
\emph{no sector structure at all}: the reset connects the whole basis into a
single enclosure. That is a real result rather than a plotting problem, and it
has a corollary for how this tier should be used. For those $150$ rules the
sector axis says nothing, and everything interesting is in the monitored map,
where the same rules spread over the whole column $a_\att\approx1$ from
$b_\att=1.0$ up to $1.89$ (right panel of Fig.~\ref{fig:f1}). The V-free
baseline behaves the other way round: it is the family that populates the
\emph{sector} plane, and it collapses almost completely in the attractor plane.

Quantitatively, with the deficit $2-a_\att b_\att$ of V4:

\begin{center}
\begin{tabular}{lrr}
\toprule
family & rules & median deficit \\
\midrule
unitary (V only) & $9$ & $0.000$ \\
V + reset & $152$ & $+0.534$ \\
V-free baseline & $78$ & $+1.000$ \\
\bottomrule
\end{tabular}
\end{center}

\noindent The unitary median is exactly zero because for a unitary rule the two
planes \emph{are} the same plane, and the median unitary rule sits on the curve.
The V-free median of exactly $+1$ says $a_\att b_\att=1$: their monitored
attractors are $O(1)$ fixed points and essentially the whole space is transient.
The V+reset family sits between the two, at $+0.53$.

\begin{figure}[htbp]
\centering
\ifgraphic{../../figures/fig_monitored_map_obc0.pdf}{0.66\textwidth}
\caption{\textbf{F2, the monitored attractor map for the quantum rules alone}
(unitary and V+reset; the V-free baseline is dropped here). MONITORED
quantities throughout. This is where the V+reset family has structure: in the
sector plane $150$ of the $160$ collapse onto $(1,2)$.}
\label{fig:f2}
\end{figure}

\begin{figure}[htbp]
\centering
\ifgraphic{../../figures/fig_sector_vs_attractor_obc0.pdf}{\textwidth}
\caption{\textbf{F3, the difference between the two scatters, drawn as a
difference.} Left: one arrow per rule, from its position in the sector plane to
its position in the monitored-attractor plane; arrow width grows with the number
of rules making the same move, and a dot means the rule does not move at all
(which is what every unitary rule does, by A2). Right: the distribution of the
deficit $2-a_\att b_\att$ by family. The V-free baseline piles up at $+1$ --- its
attractors are $O(1)$ fixed points --- while the unitary rules sit at $0$.}
\label{fig:f3}
\end{figure}

\begin{center}
\iftable{tab_r9_transient_obc0.tex}
\end{center}

\section{The corner view, and where the mass sits}
\label{sec:corner}

Two things the base plane alone cannot show. The first is the sub-leading power:
a rule with one giant sector and a rule with a \emph{linear} number of them both
sit at $a=1$ (\S\ref{sec:frontier}), and only $\alpha_{n_\wcc}$ separates them.
The second is the transient: the sector and attractor maps both plot counts and
sizes, and neither says what fraction of the basis is still moving at late times.

\subsection{Corner plots, as in R2 fig.\ 6}

R2's growth map folds the base--base plane together with the two sub-leading
powers by putting them in the margins. The main panel is base(count) against
base($D_{\max}$); the bottom margin shares its $x$-axis and plots
$\alpha_{\rm count}$ beneath it; the left margin shares its $y$-axis and plots
$\alpha_{D_{\max}}$ beside it. A rule therefore keeps its horizontal position
between the main and bottom panels and its vertical position between the main and
left panels, so one point carries the full $(b,\alpha)$ description of both
series. The same layout is worth having for both Tier-1e maps.

\begin{figure}[htbp]
\centering
\ifgraphic{../../figures/fig_sector_corner_obc0.pdf}{0.86\textwidth}
\caption{\textbf{F8, the sector map in corner form} (\rt{obc0}). Main panel: the
map of Figure~\ref{fig:f1} with the exclusion curve. Bottom margin: the sector
count's sub-leading power against its base --- the column at $a=1$ splits into
$\alpha\simeq0$ (genuinely one sector) and $\alpha\simeq1$ (the pinned-frontier
rules of \S\ref{sec:frontier}, a linear number of them), which the main panel
cannot distinguish. Left margin: the $D_{\max}$ power, where the binomial family's
$-\tfrac12$ is visible at $b=2$. Filled marker $=$ exact base.}
\label{fig:cornersec}
\end{figure}

\begin{figure}[htbp]
\centering
\ifgraphic{../../figures/fig_attractor_corner_obc0.pdf}{0.86\textwidth}
\caption{\textbf{F9, the same for the monitored attractors} (\rt{obc0}), i.e.\
terminal SCCs. \emph{Monitored only.} The dotted curve is drawn as a reference
and is \emph{not} a bound here --- terminal SCCs do not partition the basis --- and
most rules do lie below it. The V-free family (green) collapses onto
$b_{D_{\max}}=1$: their attractors are single states or short cycles however many
there are.}
\label{fig:cornerat}
\end{figure}

\subsection{Terminal SCCs against transient states}
\label{sec:rectrans}

The recurrent set $\mathrm{Rec}$ is the union of the terminal SCCs; everything
else is transient. Writing $|\mathrm{Rec}|\sim c^N$, the base $c$ answers the
question directly: $c=2$ means a finite fraction of the basis is still recurrent,
and anything below means the transient part wins exponentially.

\begin{figure}[htbp]
\centering
\ifgraphic{../../figures/fig_rec_transient_obc0.pdf}{\textwidth}
\caption{\textbf{F10, where the mass sits} (\rt{obc0}, monitored). Left: growth
base of the number of terminal SCCs against the growth base of the recurrent
\emph{mass}. Right: the same without any fit --- the recurrent fraction at the
largest $N$ computed, each family sorted.}
\label{fig:rectrans}
\end{figure}

\noindent Three readings.

\begin{itemize}
\item \textbf{The unitary rules sit exactly at $c=2$, alone.} Every state is
  recurrent, so $|\mathrm{Rec}|=2^N$ and the recurrent fraction is $1$. That is a
  consistency check rather than a finding, but it fixes the top of the plot.
\item \textbf{No dissipative rule is anywhere near it.} Not one of the $233$
  non-unitary rules with a descriptor has $c$ within $0.02$ of $2$: adding a
  single reset makes the recurrent set an exponentially vanishing share of the
  basis. The median recurrent fraction at the largest $N$ is
  $6.5\times10^{-3}$ for the V+reset family and $2.6\times10^{-4}$ for the V-free
  one.
\item \textbf{The V-free family contracts somewhat harder}, which is the opposite
  of what one might guess from the fact that it has no quantum gate, but the
  effect is modest: $51\%$ of V-free rules have a \emph{bounded} recurrent set
  ($c=1$, a fixed handful of frozen states at any $N$) against $44\%$ of the
  V+reset rules, and the two base distributions separate at $p=0.006$ on a rank
  test ($p=0.013$ on the log recurrent fraction). The Hadamard, by spreading a
  state over its sector, keeps more of the basis alive --- but by about a decade
  in the median fraction, not by an exponent.
\item \textbf{And most of that gap is one subfamily.} Splitting the V-free rules
  at whether the rule is all-reset:
  \begin{center}\small
  \begin{tabular}{lcccc}
  \toprule
  V-free subfamily & rules & bounded & median $c$ & median $|\mathrm{Rec}|/2^N$ \\
  \midrule
  all-reset (the $16$ boolean $f(\ell,r)$) & $14$ & $12$ & $1$ & $1.5\times10^{-5}$ \\
  reset $+$ identity & $64$ & $28$ & $\rho=1.3247$ & $1.3\times10^{-3}$ \\
  \bottomrule
  \end{tabular}
  \end{center}
  $1.5\times10^{-5}$ is $2^{-16}$: a \emph{single} recurrent state. An all-reset
  rule keeps no memory of the previous configuration, so it collapses to a fixed
  point. Remove those $14$ and the rest of the V-free family sits within a factor
  of $5$ of the V+reset one. The extreme contraction is a property of
  memorylessness, not of V-freeness.
\end{itemize}

\noindent \textbf{A warning about medians on this plot.} An earlier draft of the
third bullet quoted the \emph{median mass base} of each family: exactly $1$ for
V-free against $\psi=1.46557$ for V+reset. Both numbers are correct and both are
worthless. The fitted bases are not a continuous spread --- every one of them is
an algebraic number, and they pile up on a few atoms ($68$ V+reset rules at $1$,
then $21$ at $\varphi$, $21$ at $\sqrt3$, $9$ at $1.887$, $8$ at $1.848$, and
only $6$ at $\psi$). The median is a rank statistic, and rank $78$ of $155$ simply
falls inside the six-rule $\psi$ block; landing on a named constant carries no
information when every candidate value is one. It is also unstable exactly where
it was being used: the $68$ rules at base $1$ are ten short of a majority, so
moving ten rules would put the V+reset median at $1$ as well and erase the
contrast, while the V-free median of $1$ rests on $40$ of $78$ --- a majority of
one rule. Quantities with atoms of this size want a proportion or a rank test,
which is what the bullet now reports; the recurrent \emph{fraction} is continuous
and needs no such care.

\noindent A caution on what F9 and F10 are. Both are statements about the
\emph{monitored} circuit: the recurrent/transient split is a property of the
measured chain, and the unmonitored channel can keep weight on states the graph
calls transient (R5's rule 22, R7's immortal transient). The sector map F8 is the
one that is exact for both experiments.

\subsection{Weak components against terminal SCCs, directly}
\label{sec:wccscc}

F1 puts the two planes side by side and F3 draws the move between them. Neither
answers the blunter question: for one rule, how does the number of weak
components compare with the number of terminal SCCs? Two views, because they are
different claims.

\begin{itemize}
\item \textbf{Counts} (F11). The raw numbers at $N_{\max}$. Every weak component
  contains at least one terminal SCC, so $n_\scc \ge n_\wcc$ pointwise --- an
  exact inequality with no fit in it, and the figure is a check of it as much as
  a display. It holds for all $249$ rules with both descriptors.
\item \textbf{Bases} (F12). The asymptotic version. This can sit \emph{on} the
  diagonal while the counts do not: a fixed number of attractors per sector
  cancels out of the growth base.
\end{itemize}

\begin{figure}[htbp]
\centering
\ifgraphic{../../figures/fig_wcc_vs_scc_count_obc0.pdf}{\textwidth}
\caption{\textbf{F11, weak components against terminal SCCs: raw counts}
(\rt{obc0}, monitored). Left: $n_\wcc$ against $n_\scc$ at $N_{\max}$, log--log,
with the diagonal. Nothing may fall below it. Right: the ratio
$n_\scc/n_\wcc$ per rule, each family sorted --- how many attractors one sector
carries.}
\label{fig:wccsccnum}
\end{figure}

\begin{figure}[htbp]
\centering
\ifgraphic{../../figures/fig_wcc_vs_scc_base_obc0.pdf}{\textwidth}
\caption{\textbf{F12, the same comparison in growth bases} (\rt{obc0},
monitored). Rules $235$ and $249$ are off the diagonal in F11 and back on it
here, which is the distinction the two views exist to make.}
\label{fig:wccsccbase}
\end{figure}

\noindent\textbf{The sector-to-attractor map is a bijection for $237$ of $249$
rules.} That is the headline, and it is stronger than anything the base plane
shows: for $95\%$ of the rule space each weak component contains \emph{exactly
one} terminal SCC, so counting sectors and counting attractors are the same
operation and the WCC/SCC distinction adjudicated in R-T14 is vacuous. All nine
unitary rules are on the diagonal (forced: no transient, so each component is one
SCC) and so are all $80$ V-free rules.

\noindent\textbf{The twelve exceptions are all V+reset}, and they are where the
distinction earns its keep:

\begin{center}\small
\begin{tabular}{lccccl}
\toprule
rule & tuple & $n_\wcc$ & $n_\scc$ & ratio & base $n_\wcc \to$ base $n_\scc$ \\
\midrule
$203$ & \rt{VEII} & $1$ & $277$ & $277$ & $1 \to \psi$ \\
$217$ & \rt{VIEI} & $1$ & $277$ & $277$ & $1 \to \psi$ \\
$219$ & \rt{VEEI} & $1$ & $129$ & $129$ & $1 \to \psi$ \\
$36$  & \rt{IDDV} & $10$ & $595$ & $59.5$ & $1 \to \psi$ \\
$44$  & \rt{IIDV} & $17$ & $595$ & $35.0$ & $1 \to \psi$ \\
$100$ & \rt{IDIV} & $19$ & $595$ & $31.3$ & $1 \to \psi$ \\
$104$ & \rt{DIIV} & $9$ & $181$ & $20.1$ & $1 \to 1.3803$ \\
$233$ & \rt{VIIE} & $9$ & $148$ & $16.4$ & $1 \to 1.3803$ \\
$29$  & \rt{EIVD} & $22$ & $86$ & $3.9$ & $1.2147 \to \rho$ \\
$71$  & \rt{EVID} & $40$ & $86$ & $2.1$ & $1.2294 \to \rho$ \\
$235$ & \rt{VEIE} & $1$ & $2$ & $2.0$ & $1 \to 1$ \\
$249$ & \rt{VIEE} & $1$ & $2$ & $2.0$ & $1 \to 1$ \\
\bottomrule
\end{tabular}
\end{center}

\noindent Three things to read off. Rules $203$ and $217$ are the extreme case:
a \emph{single} sector --- the dynamics connects the whole basis --- carrying
$277$ distinct attractors inside it. On the sector axis they are indistinguishable
from rule $51$, which has one sector and one attractor; the monitored axis
separates them by a factor of $277$. Second, ten of the twelve are off the
diagonal in the bases too, and the jump is to a named constant in every case:
six rules go from a bounded sector count to $\psi=1.46557$ attractors, two to
$1.3803$, and $29/71$ from a fitted base to $\rho=1.32472$. The attractor count
grows exponentially inside a sector count that does not. Third, $235$ and $249$
are exactly the case F12 exists to catch: two attractors per sector at every $N$,
so the counts differ by a constant factor and the \emph{bases} agree. A base-only
comparison would have called them bijective and been wrong.

\noindent\textbf{``Twelve'' is measured at $N_{\max}$.} Both panels compare the
two counts at the largest $N$ computed, so the twelve are the rules with a
multi-attractor sector \emph{there}. That is not the same as the rules that ever
have one. Sweeping all $256$ rules at $N=11$ returns thirteen: these twelve plus
rule $37$ (\rt{EDDV}), which has the property exactly when $N\equiv2\pmod3$ and
therefore not at $N_{\max}=16$. The twelve are otherwise robust --- each has the
property at every $N$ from $6$ to $14$, bar $71$ at $N=6$. \S\ref{sec:excgraphs}
works rule $37$ through and draws it.

\subsection{Correction (2026-08-01): $\alpha$ and the base came from different
models}
\label{sec:alphafix}

Drawing these margins next to R2's exposed a defect in the fitting layer that the
base-only maps had hidden for the whole project. The descriptor reports a
\emph{pair}, $y \sim c\,N^{\alpha} b^{N}$, and the two halves were being taken
from different models. BIC fits M2, $\ln y = c + \alpha\ln N + \kappa N$, whose
$\alpha$ is whatever compensates \emph{its own} $\kappa$. Every branch that runs
afterwards then replaces the base --- with an exact integer recurrence, a parity
recurrence, the pure-exponential refit of \S\ref{sec:fitting}, or a clamp --- and
the inherited $\alpha$ was left describing a base that was no longer there.

Rule $28$'s $D_{\max}$ is the cleanest case. The series is
\[
4,\,4,\,8,\,8,\,8,\,16,\,16,\,16,\,32,\,32,\,32 \qquad (N=6\ldots16),
\]
exactly $c\cdot 2^{N/3}$ --- a pure staircase with \emph{no} prefactor. The
recurrence correctly pins the base at $2^{1/3}=1.2599$, and the descriptor then
reported $\alpha = 2.34$ alongside it, a spurious $N^{2.34}$ produced entirely by
M2's abandoned $\kappa$. Rules $39$ and $47$ carried a fictitious $\alpha=1.24$
the same way.

The fix is to hold the base that is actually being reported fixed and regress the
residual on $\ln N$ (\texttt{sectors.\_alpha\_at\_base}). Rule $28$ then gives
$\alpha = -0.03$, rule $39$ gives $0.08$, rule $47$ gives $0.03$: correctly flat,
as the series plainly are. In the sector map $111$ of $512$ $\alpha$ values move,
by a median of $0.10$ and a maximum of $1.99$.

\noindent\textbf{Nothing else moves.} No base changes, not one, and no verdict of
\S\ref{sec:hyper} changes; the defect was confined to the sub-leading power,
which is why the base-only maps of R2 and \S\ref{sec:contrast} were never wrong. Both
readings this section makes survive unchanged: the $a=1$ column still splits into
$173$ rules at $\alpha\simeq0$ and $30$ at $\alpha>0.5$, containing all eight
pinned-frontier rules, and the binomial family still carries $\alpha=-\tfrac12$ at
$b=2$ (analytic for $150/105$, and $-0.43$ from the volume-fraction discriminator
for $134$).

\noindent\textbf{R2's growth map had the same defect, and worked around it.}
Its \texttt{growth\_map.\_dissipative\_descriptor} set $\alpha=0$ by fiat for
every rule it classed exponential, on the stated grounds that ``the sub-leading
power is not identifiable under a period split''. The diagnosis was too
pessimistic --- the power is perfectly identifiable once the base is pinned; what
was unidentifiable was the \emph{stale} $\alpha$ --- but the workaround was the
right call given the value on offer, and it is why R2's fig.\ 6 margins are a
column of structural zeros for $440$ of $486$ series. With both sides repaired,
the two figures agree: the $63$ disagreements above $0.15$ fall to $2$, and those
two are rules $105$ and $150$, where R9 uses R8's derived linear sector count
($\alpha=1$ exactly) and R2 still fits $0.83$. Derivation beats fit, so R9's
value stands.

\noindent A related defect fixed at the same time: \texttt{sector\_figure}'s
\verb|__main__| guard sat in the middle of the module, before \texttt{fig\_corner}
and \texttt{fig\_recurrent\_transient} were defined, so
\verb|python -m ...sector_figure| raised \texttt{NameError} after writing F1--F6
and left F8--F10 stale on disk. Both are now pinned by tests.

\section{The survivors, drawn}
\label{sec:graphs}

Everything so far has been a growth exponent. This section draws the objects the
exponents describe: the actual transition graph of each survivor at $N=11$
($2048$ basis states), one node per state and one edge per transition.

\paragraph{Provenance.} The layout is a port of \fn{HSF/visualization\_V3.jl} and
\fn{\_V4.jl}, the Julia used for the Archipelago pictures, into
\fn{code/qca\_fragmentation/viz/dot\_sectors.py}. A port was necessary rather
than a convenience: HSF's \fn{build\_unitary\_sparse} calls
\fn{binary\_coefficients} without a mode, which restricts it to rules $0$--$15$,
and the model carries no reset operator at all --- so it cannot build a single
one of the $22$ survivors, every one of which is a V+reset rule. The conventions
do agree where they overlap, which is what makes the port faithful: HSF's
boundary gates read the missing neighbour as $\ket{0}$, our \rt{obc0}, and it
applies sites $1,3,5,\dots$ then $2,4,6,\dots$, our even-then-odd brick wall.
The files carry absolute coordinates (\fn{pos="x,y!"}) and are rendered with
\begin{center}\fn{neato -n2 -Tpng}\end{center}
where \fn{-n2} is what makes \fn{neato} honour the coordinates instead of
recomputing a layout. Intra-sector geometry is stress majorization (SMACOF),
the algorithm behind \fn{NetworkLayout}'s \fn{Stress()}.

\paragraph{Three pictures per rule.} The WCC and SCC decompositions are different
objects and this report has kept them apart throughout --- a weak component is a
\emph{sector}, exact for both experiments; a terminal SCC is a \emph{monitored
attractor}. So each rule gets three files:

\begin{itemize}
\item \fn{\_wcc} --- the whole basis, positioned and coloured by weak component.
  This is the sector decomposition of \S\ref{sec:contrast}, drawn.
\item \fn{\_scc} --- the \emph{recurrent subgraph alone}, positioned and coloured
  by terminal SCC. Transient states are omitted deliberately: they belong to no
  attractor, and for a dissipative rule they are almost the whole basis, so
  including them would bury the object of interest.
\item \fn{\_joint} --- the whole basis in the \emph{same positions} as
  \fn{\_wcc}, coloured by which attractor each state drains to, recurrent states
  solid and transient states faded. A sector holding one attractor is a single
  colour; a sector holding several is visibly striped. Black marks a state whose
  basin is shared between attractors.
\end{itemize}

\noindent The joint picture is the one that earns its place, because it makes
\S\ref{sec:wccscc}'s count concrete. Exactly six of the $22$ survivors are among
the twelve rules whose sectors carry more than one attractor --- $29$ and $71$
among the exponential ones, $44$, $100$, $104$ and $233$ among the polynomial ---
and those six are precisely the six whose joint pictures come out striped rather
than solid. The other sixteen are one colour per island.

\begin{center}
\iftable{tab_r9_graphs_obc0.tex}
\end{center}

\noindent\textbf{The same attractors, differently enclosed.} The table contains a
coincidence worth stating outright. Rules $29$ (\rt{EIVD}) and $157$ (\rt{EIVI})
differ in a single symbol --- one reset switched to the identity --- and at
$N=11$ they have \emph{literally the same $21$ attractors}, the same states in
the same classes, totalling the same $|\mathrm{Rec}|=99$. What differs is the
enclosure: $157$ puts them in $21$ sectors, one apiece, while $29$ puts them in
$14$, so seven of $157$'s sectors have merged and the merged ones hold two, three
or four attractors each. Rules $71$ (\rt{EVID}) and $199$ (\rt{EVII}) are the
identical story. Put the two \fn{\_scc} pictures side by side and they are the
same picture; put the two \fn{\_wcc} pictures side by side and they are not.
That is the WCC/SCC distinction of R-T14 in its purest form: the monitored
dynamics is unchanged and the sector structure is not, so a measurement that sees
only attractors cannot tell $29$ from $157$.

\noindent\textbf{Two more readings.} The multi-attractor rules are extreme rather
than marginal: $44$ and $100$ pack $88$ attractors into $14$ sectors with $28$ of
them inside a single sector, and all $88$ are fixed points ($|\mathrm{Rec}|=88$).
And they have large \emph{shared} basins --- $980$ of $2048$ states for $44/100$,
a full $1024$ for $233$ --- where a transient state drains to more than one
attractor. That is worth flagging against Figure~\ref{fig:f4}'s caption, which
records the shared basin as identically zero for the rules shown \emph{there};
it is zero for those rules, not in general, and the sixteen single-attractor
survivors here are zero too. A shared basin needs two attractors in one sector to
be possible at all, so the two columns switch on together.

\subsection{The twelve exceptions, drawn}
\label{sec:excgraphs}

The same three diagrams for the twelve rules of \S\ref{sec:wccscc} --- the rules
whose sectors carry more than one attractor at $N_{\max}$ --- plus rule $37$,
which has the property at $N=11$ and is explained below. Six of the twelve
($29$, $71$, $44$, $100$, $104$, $233$) are survivors and are already drawn
above, so this adds seven: $203$, $217$, $219$, $36$, $235$, $249$, $37$.

\begin{center}
\iftable{tab_r9_exceptions_obc0.tex}
\end{center}

\noindent\textbf{Read the $N$.} \S\ref{sec:wccscc}'s table quotes the largest $N$
computed, $N_{\max}=16$, where $203$ and $217$ carry $277$ attractors. This table
is at $N=11$, where they carry $41$. The two are not in conflict; the attractor
count is growing like $\psi^N$ and the drawings are of the smaller system.

\noindent\textbf{For five of the six, the sector picture is a single blob.}
Rules $203$, $217$, $219$, $235$ and $249$ have $n_\wcc=1$ at $N=11$: the
dynamics connects all $2048$ basis states into one weak component, so the
\fn{\_wcc} file is one undifferentiated island and carries no information at all.
Everything is in the other two. This is the most direct illustration in the
report of \S\ref{sec:pileup}'s point that the sector axis is uninformative for
the V+reset family --- here it is not merely coarse, it is a constant.

\noindent\textbf{What dominates $203$ and $217$ is the shared basin.} Their
joint picture is a single island with $41$ small colour cores scattered through
it --- and more than half of it, $1184$ of the $2048$ states, carrying the
shared-basin colour, black faded to grey because those states are transient. So
the majority of the basis is not merely transient;
it is transient to \emph{no particular attractor}, able to reach several
depending on the measurement record. Only $50$ states are recurrent. A rule whose
sector map reports ``one sector, nothing to see'' turns out to have most of its
state space in ambiguous drainage.

\noindent\textbf{$235$ and $249$ are the cleanest picture in the set.} One
sector, \emph{two} attractors, both fixed points, and all $2048$ states draining
into them with $768$ of them able to reach either. The \fn{\_wcc} picture is one
undifferentiated blob; the \fn{\_joint} picture is that same blob resolved into
one basin left, the other right, and a dark central third --- not a seam but
$37.5\%$ of the basis --- that belongs to neither. That is the entire content of
\S\ref{sec:wccscc}'s
count-versus-base distinction in one image: the counts differ by a factor of two
at every $N$, which is why F11 puts these rules off the diagonal, and a constant
factor cancels out of a growth base, which is why F12 puts them back on it.

\noindent\textbf{$36$ is the extreme of concentration.} Ten sectors, but one of
them holds $1999$ of the $2048$ states and $79$ of the $88$ attractors. Its joint
picture is a single giant island shot through with $79$ colours and nine crumbs
around it. All $88$ attractors are fixed points ($|\mathrm{Rec}|=88$).

\noindent\textbf{Equal rows are not always twin rules.} Six of the twelve appear
as identical-looking pairs, and they are identical to different depths, so the
table alone does not settle it. Rules $44$ (\rt{IIDV}) and $100$ (\rt{IDIV}) do
share an attractor set exactly --- the same $88$ fixed points --- with the same
sector \emph{sizes}, but their sector partitions and their transition graphs
differ, and so do their shared basins despite both containing $980$ states.
Rules $203$/$217$ and $235$/$249$ agree in every count in the table and do
\emph{not} share an attractor set at all. Only the $29$/$157$ and $71$/$199$
pairs of \S\ref{sec:graphs} have the strong property, and note the complement
there: $157$ has no shared basin whatever ($0$ states), because with one
attractor per sector a shared basin is impossible, while $29$ has $128$.

\noindent\textbf{A thirteenth rule, and what ``twelve'' actually means.}
Checking the drawings against the rule space turned up a rule the count in
\S\ref{sec:wccscc} does not contain. Sweeping all $256$ rules at $N=11$ for a
sector holding more than one terminal SCC returns \emph{thirteen}: the twelve,
and rule $37$ (\rt{EDDV}). It is not an error in either place. The twelve are
read off the descriptor comparison at $N_{\max}=16$, and $37$ has the property
only at $N\equiv2\pmod 3$:
\begin{center}\small
\begin{tabular}{lccccccccccc}
\toprule
$N$ & $6$ & $7$ & $8$ & $9$ & $10$ & $11$ & $12$ & $13$ & $14$ & $15$ & $16$ \\
\midrule
$n_\scc$ of $37$ & $1$ & $1$ & $2$ & $1$ & $1$ & $2$ & $1$ & $1$ & $2$ & $1$ & $1$ \\
\bottomrule
\end{tabular}
\end{center}
Since $16\equiv1\pmod 3$, rule $37$ has one attractor per sector exactly at the
size the comparison is made, and is invisible to it. It has one sector at every
$N$, so at $N=8,11,14$ it is a one-sector two-attractor rule --- the same shape
as $235$ and $249$ --- and at every other $N$ it is unremarkable. Its picture is
the extreme of the whole set: two fixed points, and $1920$ of the $2048$ states
--- $94\%$ of the basis --- in the shared basin, so almost nothing in the space
is committed to one attractor rather than the other.

So \textbf{``the twelve'' is a statement at $N_{\max}$, not a theorem about the
rule space}, and the honest form of \S\ref{sec:wccscc} is that twelve rules have
a multi-attractor sector \emph{at the size where the growth comparison is made}.
The twelve themselves are robust --- every one has the property at every $N$ from
$6$ to $14$, with the single exception of $71$ at $N=6$ --- and $37$ is the only
rule found that oscillates. A sweep that sampled $N\equiv2\pmod 3$ instead would
have reported thirteen. Rule $37$ is drawn here with the others.

\noindent A caveat on reading these by eye. The joint colouring gives one hue per
attractor, which works while there are tens of them; at $41$ ($203$/$217$) the
colours are still separable, but the palette is $277$-wide at $N_{\max}$ and its
closest pair then differs by an RGB distance of $2$, which is no distance at all.
So the joint picture is evidence that a sector is shattered into many basins, and
the table is the place to read \emph{how many}. It was never going to be possible
to distinguish $277$ colours in one image.

\subsection{Files and regeneration}

The renderings are not embedded --- $84$ pictures at this size would dominate the
document --- and are not tracked either: the directory is $152$\,MB of \fn{.dot}
and \fn{.png} that no reader needs in the history. They live in
\fn{figures/sector\_graphs/} as
\fn{N11\_\{rule\}\_V4\_obc0\_\{wcc,scc,joint\}.\{dot,png\}}, and the placement
RNG is seeded per rule, so
\begin{center}
\fn{python -m qca\_fragmentation.viz.dot\_sectors -{}-N 11 -{}-bc obc0 -{}-rules all}
\end{center}
reproduces all $84$ byte-for-byte in about twenty minutes. \fn{-{}-rules all} is
the union of the two groups and is $28$ rules, not $34$, because the six above
are in both. \fn{-{}-rules survivors} and \fn{-{}-rules exceptions} draw one
group each, \fn{-{}-variant V3} gives the core-periphery layout instead of the
archipelago one, and \fn{-{}-rules 29,157} draws a chosen pair. (The directory
also holds a V3 set for the $29$/$157$ and $71$/$199$ twins, hence $96$ files
rather than $84$.)

\section{Clustering the dissipative rules by their coherent correspondent}
\label{sec:coherent}

Every V+reset rule has a \emph{coherent part}. A symbol-table entry says what
happens at a site for one neighbour pattern: \rt{I} does nothing, \rt{V} applies
the Hadamard, \rt{D} resets to $\ket0$, \rt{E} resets to $\ket1$. Switching the
dissipation off means the reset entries stop acting,
\begin{equation}
\rt{D},\ \rt{E}\;\longmapsto\;\rt{I},
\label{eq:coherentpart}
\end{equation}
which leaves only \rt{I} and \rt{V} --- always one of the sixteen unitary rules.
So \eqref{eq:coherentpart} assigns each dissipative rule a unitary
\emph{correspondent}, and the assignment partitions the rule space exactly: a
unitary parent with $v$ Hadamard entries has
\begin{equation}
3^{4-v}-1
\end{equation}
V+reset children, each non-\rt{V} entry independently becoming \rt{I}, \rt{D} or
\rt{E}, minus the all-\rt{I} case which is the parent itself. Summing over the
fifteen unitary rules with $v\ge1$ gives $4\cdot26+6\cdot8+4\cdot2+0=160$, the
whole V+reset family, in $14$ non-empty clusters
($\rt{VVVV}$ has no \rt{I} to spoil, so it has no children). This is verified
against the rule tables rather than assumed.

\subsection{What adding dissipation does}

\begin{center}
\iftable{tab_r9_dissip_obc0.tex}
\end{center}

\noindent The table splits the fourteen clusters cleanly in two.

\begin{itemize}
\item \textbf{Ten parents already sit at $(1,2)$} --- one sector holding the
  whole space --- so there is nothing for a reset to destroy, and all $56$ of
  their children stay exactly where the parent is.
\item \textbf{Four parents have genuine sector structure} --- the $v=1$ rules
  $156$, $198$, $108$, $201$ --- and they hold $104$ children between them.
  \emph{Not one of the $104$ keeps its parent's position.} $94$ collapse all the
  way to $(1,2)$; the remaining $10$ move to some intermediate point but never
  stay --- and eight of those ten are still \emph{exponentially} fragmented,
  which is \S\ref{sec:openfrag}.
\end{itemize}

\noindent That is the sharpest statement this tier produces about dissipation:
\textbf{adding any reset to a rule that has sector structure destroys that
structure, without exception over the $104$ cases available}. It also explains
the pile-up of \S\ref{sec:pileup} arithmetically --- the $150$ V+reset rules at
$(1,2)$ are exactly the $94$ that collapsed plus the $56$ that were never
anywhere else.

\subsection{Fragmentation in an open quantum system}
\label{sec:openfrag}

The $104$ children of the four structured parents split $94$ / $10$: ninety-four
collapse to $(1,2)$, and ten land somewhere in between. Those ten are the
interesting ones, and eight of them are interesting for a specific reason ---
their sector count still grows \emph{exponentially}.

\begin{center}
\iftable{tab_r9_openfrag_obc0.tex}
\end{center}

\noindent Why this is worth stating as a result rather than a curiosity. A sector
is an \emph{enclosure}: $\operatorname{span}(\text{WCC})$ is invariant under every
Kraus operator, so an exponentially growing sector count is a statement about the
unmonitored channel just as much as about the measured circuit
(\S\ref{sec:tiers}). These eight rules therefore exhibit \textbf{Hilbert-space
fragmentation as open quantum systems} --- not a unitary circuit that someone
later added noise to, and not an artefact of conditioning on measurement records.
Each contains a genuine reset, and each has a coherent correspondent that is
itself fragmented, so the fragmentation is inherited rather than accidental.

The structure of the list is tidier than the selection deserves:

\begin{itemize}
\item They come from all \textbf{four} of the structured parents --- the $v=1$
  rules $156$, $198$, $108$, $201$ --- three, three, one and one respectively.
  No $v\ge2$ parent contributes, because none of them has sector structure to
  pass on in the first place.
\item The bases are \textbf{exact algebraic numbers}, from integer recurrences,
  for six of the eight: the \textbf{plastic number} $\rho=1.32472$, the real root
  of $x^3=x+1$, for $28$, $70$, $157$, $199$; and the \textbf{supergolden ratio}
  $\psi=1.46557$, the real root of $x^3=x^2+1$, for $73$ and $109$. For the four
  $\rho$ rules the reset \emph{raises the degree} of the characteristic
  polynomial: their parents $156$ and $198$ count sectors by
  $x^2=x+1$ and the children count by $x^3=x+1$. That pattern does \emph{not}
  extend to $73$ and $109$, whose parents $201$ and $108$ already count by a
  cubic, $x^3=2x^2-x+1$; there the degree is unchanged and only the polynomial
  moves, to $x^3=x^2+1$. So there is no single statement covering all eight ---
  see \S\ref{sec:whichconstant}.
\item They pair up. $\{29,71\}$ share $\approx(1.22,1.81)$, and $73$, $109$ are
  each their own reflection partner. Within a pair the two rules are reflections
  of one another, one built on \rt{D} and the other on \rt{E}. On the $b$ axis
  the four $\rho$ rules collapse onto a \emph{single} value: $70$ and $199$ ---
  and their reflections $78$ and $157$ --- all have $D_{\max}=2\cdot3^{N/2-1}$ on
  even $N$, i.e.\ $b=\sqrt3$ exactly.
\item Rule $28$ (and its reflection $92$) is the one case in the list where the
  two parities do not share a base. Its $D_{\max}$ is
  $14,28,38,76,108,208,324,568,972,1552,2916$; the even branch settles to a ratio
  of exactly $3$ per two sites, so $b_{\rm even}=\sqrt3$ like the others, while
  the odd branch settles to $1+\sqrt3$ per two sites, i.e.\
  $b_{\rm odd}=\sqrt{1+\sqrt3}=1.6529$. The single number in the table,
  $1.68818$, is a fit through both branches and should be read as ``between the
  two''. Rule $28$ is also the one whose $\varphi$ prediction \S\ref{sec:r28}
  falsified --- and the plastic number that replaced it is the same constant that
  appears here for three other rules.
\end{itemize}

\paragraph{Correction, 2026-07-31.} An earlier version of this section reported
$b=2$ \emph{exactly} for $28$ and $70$, and $b=2$ for $132$, $133$, $164$, $222$,
$78$ and $92$ as well: eight $D_{\max}$ descriptors in all. That came from the
volume-fraction discriminator (\S\ref{sec:fitting}), which decides whether
$D_{\max}=c\,2^N N^\alpha$ by asking whether $D_{\max}/2^N$ is better explained
by a power of $N$ than by an exponential, and which was being run on a six-point
tail. Over so short a window $\ln N$ is nearly affine in $N$ and the power-law
model can win against an exponential decay that is genuinely there. Requiring it
to win on the \emph{full} series as well removes all eight, and the corrected
numbers are the ones tabulated above. The direction of the error was uniform:
each of the eight has a $D_{\max}$ that is an exponentially shrinking fraction of
the space --- $W_{222}$ falls from $0.328$ of the basis at $N=6$ to $0.010$ at
$N=16$ --- and none was a fixed fraction. No verdict changes: all eight were
already \emph{above} the exclusion curve and remain so, with products now
$2.20$--$2.33$ instead of $2.65$--$3.24$. The guard was found while applying the
same fitting layer to the X-gate cycle series in R10~\S3, where its failure mode
is far more visible.

\noindent The remaining two of the ten, $6$ (\rt{IVDD}) and $20$ (\rt{IDVD}),
have $a=1$: their sector counts are linear ($4,5,5,6,6,7,7,8,8,9,9$), so they keep
a little structure but are not fragmented.

\paragraph{A caution on the two fitted rows.} For $29$ and $71$ the base comes
from a parity-resolved fit, not a recurrence, so $1.2147$ and $1.2294$ should be
read as ``about $1.22$'' and not as two distinct constants; their series
oscillate strongly with $N$ ($3,6,5,9,7,14,10,21,15,31,22$ for rule $29$) and each
parity branch has only five or six points. That they are exponential is solid ---
both branches grow geometrically --- but the value is not pinned, and if they
follow the pattern of the other six the honest guess is that both are the same
algebraic number.

\subsection{Which constant describes what: a naming correction}
\label{sec:whichconstant}

Earlier drafts of this report called the four structured parents the
``$\varphi$-fragmented'' rules. That is wrong, and worth correcting explicitly
because the error is easy to make: $\varphi$ describes \emph{different observables}
in the two pairs, and no single constant describes a rule's fragmentation at all.
The exact recurrences:

\begin{center}
\begin{tabular}{rlll}
\toprule
rule & sector count $n_\wcc$ & largest sector $D_{\max}$ \\
\midrule
$156$, $198$ & $21,34,55,89,\ldots$ & $6,9,12,16,20,27,\ldots$ \\
 & Fibonacci, $a_n=a_{n-1}+a_{n-2}$ & \emph{no} integer recurrence \\
 & base $\varphi$, root of $x^2=x+1$ & base $4^{1/5}$, from theory (R2~\S3) \\
\midrule
$108$, $201$ & $37,65,114,200,\ldots$ & $8,13,21,34,55,\ldots$ \\
 & $a_n=2a_{n-1}-a_{n-2}+a_{n-3}$ & Fibonacci, $a_n=a_{n-1}+a_{n-2}$ \\
 & base $1.75488$, root of $x^3=2x^2-x+1$ & base $\varphi$ \\
\bottomrule
\end{tabular}
\end{center}

\noindent Two separate problems with the old label.

\begin{itemize}
\item \textbf{Across rules, $\varphi$ changes role.} For $156$ and $198$ it is the
  base of the sector \emph{count} --- there are Fibonacci-many sectors. For $108$
  and $201$ it is the base of the largest sector \emph{size} --- the biggest
  sector is Fibonacci-sized --- while their sector count grows as the root of
  $x^3=2x^2-x+1$. Calling both pairs ``$\varphi$-fragmented'' asserts a shared
  property they do not have.
\item \textbf{Within one rule, the two quantities have unrelated mechanisms.}
  Rule $156$'s sector count is a Fibonacci word count; its largest sector is
  \emph{not} an integer recurrence at all, and its base comes from the
  room-packing saddle point of R2~\S3, an optimisation over room length that has
  nothing to do with $\varphi$. Naming the rule after one constant hides the fact
  that its two structural laws are different laws.
\end{itemize}

\noindent So the group is referred to below as \textbf{the four exponentially
fragmented parents}, and any base is quoted together with the observable it
belongs to. The pairs are, in the notation $(a,b)$ of \S\ref{sec:hyper},
\[
156,198:\ (\varphi,\ 4^{1/5}),\ ab=2.1350;\qquad
108,201:\ (1.75488,\ \varphi),\ ab=2.8394 ,
\]
so both sit above the exclusion curve, but for structurally different reasons ---
the first pair by having many sectors, the second by having large ones. That
$\varphi$ appears in both lines is a coincidence of which observable one looks at,
not evidence of a common mechanism.

\subsection{The base plane hides a whole tier, and 150/105 are in it}
\label{sec:polytier}

The bookkeeping above is done in the base plane, and the base plane cannot tell a
linear sector count from a single sector: both give $a=1$. So rule $150$ ---
which has $\lfloor(N+1)/2\rfloor+1$ sectors, nine of them at $N=16$, the
domain-wall charge of R8 --- is plotted on top of rule $51$, which has exactly
one. Calling both of them ``already at $(1,2)$, nothing to destroy'' is wrong for
$150$. Survival has to be asked in terms of the growth \emph{class}.

\begin{center}
\iftable{tab_r9_survclass_obc0.tex}
\end{center}

\noindent Three quite different fates, which the $(1,2)$ cell had conflated:

\begin{itemize}
\item \textbf{Exponential parents} ($108$, $156$, $198$, $201$), $104$ children:
  $8$ stay exponential (\S\ref{sec:openfrag}), \textbf{$14$ degrade to polynomial
  but keep growing}, $82$ collapse to a constant number of sectors.
\item \textbf{Polynomial parents} --- $150$ and $105$ with the wall charge
  ($\alpha=1$), $60$ and $102$ with the frontier charge of \S\ref{sec:frontier}
  (fitted $\alpha\approx0.82$, same dyadic tower) --- $32$ children:
  \textbf{$0$ survive.} Every one drops to a
  \emph{constant} sector count. For $150$ specifically, six of its eight children
  fall to a single sector and the other two ($22$~\rt{IVVD}, $182$~\rt{IVVE}) to
  two sectors; all eight of $105$'s fall to one. The domain-wall charge does not
  tolerate a reset anywhere.
\item \textbf{Constant parents} (six rules), $24$ children: nothing to lose ---
  though note even rule $54$'s two sectors become one.
\end{itemize}

\noindent So the answer for $150$ and $105$ is blunt: \textbf{nothing survives}.
That is a \emph{stronger} statement than for the exponential parents, where
$8+14$ of $104$ children keep some growth. Polynomial sector structure is more
fragile than exponential sector structure, not less.

\paragraph{But the wall charge is a sink, not a source.} The fourteen polynomial
survivors are worth looking at, because six of them have \emph{exactly} rule
$150$'s sector series $4,5,5,6,6,7,7,8,8,9,9$ and two have rule $105$'s:

\begin{center}
\iftable{tab_r9_polysurv_obc0.tex}
\end{center}

\noindent Those six are $20$, $148$, $158$ (from $156$) and $6$, $134$, $214$
(from $198$) --- and four of them are exactly the mixed rules of
\S\ref{sec:r150family} that share $150$'s whole multiset. So dissipation can
degrade the exponential fragmentation of $156$/$198$ \emph{down to} the
rule-$150$ wall charge, while a
reset applied \emph{to} the wall charge destroys it outright. The wall charge sits
one level below the exponential rules in a hierarchy
\[
\text{exponential count }(\varphi\ \text{or}\ 1.75488)\;\longrightarrow\;
\underbrace{\rho,\ \psi}_{8}\;\longrightarrow\;
\underbrace{\text{wall charge}}_{14}\;\longrightarrow\;
\underbrace{\text{one sector}}_{82},
\]
and it is a terminal rung: reachable from above, but with no dissipative
descendants of its own.

\subsection{A second polynomial mechanism: the pinned frontier}
\label{sec:frontier}

Six of the fourteen polynomial survivors --- $44$, $100$, $110$, $124$, $188$,
$230$ --- are \emph{not} degraded copies of the rule-$150$ wall charge. They carry
a different conserved quantity, and two \emph{unitary} rules, $60=\rt{IIVV}$ and
$102=\rt{IVIV}$, carry the same one, so it is not produced by dissipation at all.
The sector sizes give it away immediately.
At $N=12$, rule $110$ has thirteen sectors of sizes
\[
2048,\,1024,\,512,\,256,\,128,\,64,\,32,\,16,\,8,\,4,\,2,\,1,\,1 ,
\]
a \textbf{dyadic tower} rather than a binomial row. The charge is the
\textbf{position of the extremal excitation}: verified at $N=9$, every sector is a
level set of the index of the rightmost $1$ for $110$, $230$ and $102$, and of the
leftmost $1$ for $124$, $188$, $44$ and $60$. The outermost excitation is pinned by the
reset and everything inside it is free, which gives
\begin{equation}
n_\wcc = N+1,\qquad D_{\max}=2^{N-1}\ \text{exactly},\qquad
\frac{D_{\max}}{2^N}=\tfrac12 .
\end{equation}

\begin{center}
\iftable{tab_r9_frontier_obc0.tex}
\end{center}

\noindent Set against the wall charge, the contrast is sharp, and both are
invisible in the base plane because both have $a=1$, $b=2$:

\begin{center}
\begin{tabular}{lll}
\toprule
 & rule-$150$ family & frontier family \\
\midrule
conserved quantity & domain-wall number & position of the extremal excitation \\
$n_\wcc$ & $\lfloor(N+1)/2\rfloor+1\sim N/2$ & $N+1$ \\
sector sizes & binomial $\binom{N+1}{w}$, even $w$ & dyadic $2^{N-1},\ldots,2,1,1$ \\
$D_{\max}$ & $\binom{N+1}{\sim N/2}\sim2^N/\sqrt N$ & $2^{N-1}$ exactly \\
$D_{\max}/2^N$ & $\sim N^{-1/2}\to0$ & $\tfrac12$, constant \\
\bottomrule
\end{tabular}
\end{center}

\noindent The frontier family has \emph{twice} as many sectors asymptotically and
a largest sector that stays a fixed half of the space rather than shrinking. Two
of the six dissipative members are less clean: $44$ splits each frontier class
further at odd $N$ ($n_\wcc=N+1$ for even $N$, $N+3$ for odd), and $100$ has $N+3$
sectors with a short non-dyadic tail ($\ldots,26,16,6,6,4,2$), so its charge is
the leftmost excitation plus a refinement.

\paragraph{Two of them are unitary, which reclassifies them.} $60$ and $102$ were
listed in \S\ref{sec:polytier} among the polynomial parents beside $150$ and
$105$, on the strength of a fitted $\alpha\approx0.82$. That grouping is wrong:
their sector sizes are the same dyadic tower and their charge is the same extremal
excitation, so they belong here, and the differing $\alpha$ only reflects the
fitter seeing one series through a different window. It also means the frontier
charge exists already in the \emph{unitary} rule space --- it is a destination
reachable both directly ($60$, $102$) and as a dissipative degradation of
the exponential rules ($110$, $124$, $188$, $230$ from $108$, $156$, $198$) ---
and not something a reset creates.

Since sectors are enclosures, this is an open-system statement in the sense that
matters here: the pinned frontier is a conserved quantity of the \emph{channel},
not of a measurement record. It is a genuinely different route to polynomial
fragmentation than the wall charge, and both are reachable from the same
exponentially fragmented parents --- $108$ contributes four frontier rules while
$156$ and $198$ contribute the wall-charge ones.

\paragraph{A caveat this report cannot settle.} The charge here is the position of
the \emph{extremal} excitation, and that is a quantity only an open chain has ---
a ring has no leftmost site. So the frontier structure is, on its face, tied to
the boundary in a way the domain-wall number of \S\ref{sec:r150family} is not:
the wall number is defined by $b_j=x_j\oplus x_{j+1}$ everywhere in the bulk,
whereas ``the first $1$ from the left'' presupposes a left. Whether the frontier
sectors are bulk structure or an artefact of \rt{obc0} is therefore a real
question, and it is not one this report can answer, because answering it means
changing the boundary condition. It is deferred to the \rt{pbc} report, and until
then \S\ref{sec:frontier} should be read as a statement about open chains only.
The same question does not arise for the eight exponentially fragmented rules of
\S\ref{sec:openfrag}, whose charges are not defined by reference to an end, but
confirming that is also \rt{pbc} work.


\paragraph{And 105 is not 150's partner.} R8 treats them as reflection partners
and this report inherited that, but the sector multisets differ at
$N\equiv1\pmod4$. Rule $105=\rt{VIIV}$ fires when the two neighbours
\emph{agree}, where $150=\rt{IVVI}$ fires when they differ, and at those $N$ the
sectors of $105$ carry \emph{odd} wall number: at $N=9$ its multiset is
$[252,120,120,10,10]=\binom{10}{\text{odd }w}$ against $150$'s
$[210,210,45,45,1,1]=\binom{10}{\text{even }w}$, and likewise at $N=13$. At all
other $N$ in range the two agree exactly. The growth is linear either way, which
is all the base captures --- another thing the base plane hides.

\begin{figure}[htbp]
\centering
\ifgraphic{../../figures/fig_dissip_sector_obc0.pdf}{\textwidth}
\caption{\textbf{F6, the move in the SECTOR plane, clustered by coherent
correspondent.} One panel per unitary rule: the star is the coherent rule
(all resets switched off, \eqref{eq:coherentpart}), the dots are its V+reset
children, and each arrow is the move that turning the resets on produces. Dot
area and arrow width grow with the number of children making the same move.
Top row and the two leftmost middle panels: parents with sector structure, every
arrow pointing away and up-left towards $(1,2)$. Remaining panels: parents
already at $(1,2)$, where the children coincide with the parent and no arrow is
drawn. Gold rings mark the eight children that remain exponentially fragmented
(\S\ref{sec:openfrag}). These coordinates are exact for both experiments.}
\label{fig:f6}
\end{figure}

\begin{figure}[htbp]
\centering
\ifgraphic{../../figures/fig_dissip_attractor_obc0.pdf}{\textwidth}
\caption{\textbf{F7, the same clusters in the MONITORED attractor plane.} Here
the moves are large for every cluster, including the eight whose sector
coordinates do not move at all --- which is the point of keeping both planes:
dissipation is invisible on the sector axis for those rules and plainly visible
on the attractor axis. Six panels have no star: those parents
($54$, $57$, $99$, $147$, $153$, $195$) are ergodic at \rt{obc0}, so Tier 1a
never produced an attractor descriptor for them and the cluster has no origin to
draw from. Monitored quantities throughout.}
\label{fig:f7}
\end{figure}

\section{Rule 150's wall charge survives six dissipative rules}
\label{sec:r150family}

The sector axis makes a comparison possible that the attractor axis could not,
and it turns up an exact coincidence. Consider the nine rules of the form
\rt{I}$\,\cdot\,\cdot\,$\rt{I}, i.e.\ those where a site is inert unless its two
neighbours differ. Seven of them --- including six \emph{dissipative} ones ---
have \textbf{exactly} rule $150$'s sector-size multiset, at every
$N=6\ldots16$:

\begin{center}
\begin{tabular}{llc}
\toprule
rule & tuple & sector multiset $=$ that of $W_{150}$ \\
\midrule
$134$ & \rt{IVDI} & \cmark \\
$142$ & \rt{IEDI} & \cmark \\
$148$ & \rt{IDVI} & \cmark \\
$150$ & \rt{IVVI} & --- (unitary reference) \\
$158$ & \rt{IEVI} & \cmark \\
$212$ & \rt{IDEI} & \cmark \\
$214$ & \rt{IVEI} & \cmark \\
$132$ & \rt{IDDI} & \xmark\ ($n_\wcc=2584$, $D_{\max}=322$ at $N=16$) \\
$222$ & \rt{IEEI} & \xmark\ ($n_\wcc=988$, $D_{\max}=686$ at $N=16$) \\
\bottomrule
\end{tabular}
\end{center}

\noindent So the domain-wall charge of R8 --- $|S_w|=\binom{N+1}{w}$ on even $w$,
$\lfloor(N+1)/2\rfloor+1$ sectors --- is not a property of unitarity at all. It
survives replacing \emph{either or both} Hadamards by a reset, and the two
exceptions are precisely the rules whose two active symbols are the \emph{same}
reset type (\rt{DD} or \rt{EE}). This is the sector-level counterpart of the
Tier-2 observation (R-T13) that $W_{148}$ and $W_{134}$ inherit an exact
decoherence-free subspace of dimension $\lfloor(N+1)/2\rfloor+1$, which is
exactly C150's sector count: here the whole multiset, not just its cardinality,
is shown to agree, and for six rules rather than two.

\paragraph{A conjecture that fails.} The obvious generalisation --- that the
sector partition depends only on which of the four symbols are \rt{I}, since
\rt{V}, \rt{D} and \rt{E} all generate the same \emph{undirected} flip edge
$x\leftrightarrow x\oplus2^i$ --- is \textbf{false}. Grouping all $256$ rules by
their \rt{I}-pattern gives $16$ groups, and only three collapse to a single
sector signature: $\{140,156,220\}$, $\{196,198,206\}$ and the singleton
$\{204\}$. (Those first two are exactly the six rules of
\S\ref{sec:inconclusive} whose \emph{sector count} has base $\varphi$.) The reason the
argument fails is the brick-wall composition: the odd-layer symbol is read off
the \emph{post-even-layer} state, and a \rt{V} leaves a branch supported on both
values of the flipped bit where a \rt{D} or \rt{E} leaves only one, so the
downstream firing pattern differs even though the single-site undirected edge
does not.

\section{Growth classes and exact bases}
\label{sec:growth}

\begin{center}
\iftable{tab_r9_classes_obc0.tex}
\end{center}

\begin{center}
\iftable{tab_r9_exact_obc0.tex}
\end{center}

\section{Basins, and a gap in the Tier-1a archive}
\label{sec:basins}

Basins are a \textbf{monitored} observable throughout this section.

The sum rule $\sum_k|\text{basin}_k|+|\text{shared}|=2^N$ turned out to be
\emph{unanswerable} for $76$ \rt{obc0} units. The cause is a schema gap rather
than a physics failure: Tier 1a caps \fn{sizes\_basins} at $2048$ entries and,
unlike \fn{sizes\_recurrent}, never stored a basin histogram, so for a unit
whose basin list hit the cap the multiset is genuinely lost. Every one of the
$76$ failures is a truncated record and no untruncated record fails.

Asserting the check away would have hidden a real gap. Instead those units were
recomputed from the rule definition and stored alongside, in
\fn{results\_basins/}, this time \emph{with} the histogram, and cross-checked
against the Tier-1a fields that did survive (\fn{n\_scc}, \fn{n\_recurrent},
\fn{shared\_basin\_size}, \fn{transient\_depth}).

\begin{center}
\iftable{tab_r9_basins_obc0.tex}
\end{center}

\noindent \textbf{The gap is closed}: no unit remains truncated-and-unrecomputed.
The checks that are still unavailable are unavailable for reasons that are not
defects. Fifty-six of them are units where the Tier-1a Tarjan pass took its
ergodic early exit and so never computed basins at all --- there the check is
\emph{undefined}, not missing. The remaining twenty-five are the $N=17$--$18$
units of the targeted extension, which lie beyond where the Tier-1a sweep ever
ran; the sector quantities there are complete, only the monitored basin
comparison has nothing to compare against.

\begin{figure}[htbp]
\centering
\ifgraphic{../../figures/fig_basins_obc0.pdf}{\textwidth}
\caption{\textbf{F4, basin structure (MONITORED).} Largest basin, shared basin
and transient fraction against $N$ for selected rules. Two things to read off.
The shared basin is \emph{identically zero} for every rule shown: each transient
state funnels to a unique terminal SCC, so the basins alone already partition the
transient part and no state is ambiguous between attractors. And the two unitary
rules ($156$, $201$) sit at transient fraction $0$ by construction --- every state
is recurrent --- which is the cleanest illustration of why the transient fraction
is a monitored quantity and not a sector one. Tier 1d's certificate T3 gives
certified zeros for the \emph{unmonitored} basins of the V-free rules, so for
those $80$ rules the monitored and channel populations coincide and these curves
are channel-exact; for the rest they are monitored-only.}
\label{fig:f4}
\end{figure}

\section{Validation}
\label{sec:valid}

\begin{center}
\iftable{tab_r9_validation_obc0.tex}
\end{center}

\begin{figure}[htbp]
\centering
\ifgraphic{../../figures/fig_validation_obc0.pdf}{\textwidth}
\caption{\textbf{F5, validation panel.} Left: distribution of the hyperbola
margin $ab-2$. Right: pass rates for the sum rule A1 (residual identically
zero), the unitary agreement A2, and the monitored basin sum rule.}
\label{fig:f5}
\end{figure}

\noindent On the cross-tier assertions: A2 (for unitary rules, $n_\wcc=
n_{\rm recurrent}=n_{\rm scc}$ with identical size multisets) holds wherever
both records exist. A3 ($n_\wcc\le n_{\rm scc}$) is structural and holds
throughout. A4 is asserted only in the direction that is actually true --- every
sector contains at least one terminal SCC, hence $n_{\rm recurrent}\ge n_\wcc$;
the reverse bound is false in general and is not asserted.

\section{Comparison hooks for Tier 2}
\label{sec:hooks}

The numbers the wall grammar has to reproduce, in order of how sharply they
discriminate:

\begin{enumerate}
\item \textbf{Rule 28's plastic number.} Equation~\eqref{eq:r28} and
  $\rho=1.32472$, exact. The transparency argument that predicted $\varphi$ is
  wrong; a correct grammar must explain why one reset symbol takes
  $\varphi\to\rho$, i.e.\ why $x^2=x+1$ becomes $x^3=x+1$.
\item \textbf{The six rules with a Fibonacci sector count.} $140$, $156$, $196$,
  $198$, $206$, $220$ all have $a=\varphi$ exactly at \rt{obc0}; the grammar should say why these six and
  no others, and should predict $b$, which the hyperbola confines to
  $[1.2361,1.2369]$.
\item \textbf{The exact-base census} of \S\ref{sec:growth}: every rule whose
  sector count satisfies an integer recurrence, with the recurrence.
\item \textbf{The on-curve population.} $189$ of $256$ rules sit at $ab=2$
  exactly, i.e.\ their sector sizes are all comparable. Which symbol tables force
  that is a grammar question.
\item \textbf{Rule 150's family} (\S\ref{sec:r150family}). The wall charge is
  shared by six dissipative rules and broken by exactly two, \rt{IDDI} and
  \rt{IEEI}. A grammar that produces the domain-wall counting should say why the
  same-reset-type pair is the one that breaks it.
\item \textbf{Which V+reset rules escape the collapse} (\S\ref{sec:pileup}).
  Only $10$ of $160$ keep any sector structure and only $8$ stay
  exponentially fragmented (\S\ref{sec:openfrag}); a grammar should predict that
  list, and it is a short one.
\end{enumerate}

\section{Open items}

\begin{enumerate}
\item \textbf{\rt{pbc}.} This report is \rt{obc0} only, by instruction. The ring
  is a separate report; note that rule $156$'s Lucas law and rule $22$'s
  attractor count are \rt{pbc} statements already verified here in passing.
\item \textbf{Separating $3^{1/4}$ from $4^{1/5}$ empirically} for the $\varphi$
  family. R2~\S3 settles it by derivation and this report follows that; the two
  candidates differ by $0.26\%$, so confirming it from data alone would need far
  more than the eleven sizes available here.
\item \textbf{$W_{90}$ and $W_{165}$} (both \rt{DEED}/\rt{EDDE}) have irregular
  sector counts that no growth law fits. Their structure is presumably
  number-theoretic in $N$; nothing in the engine is wrong (A1 holds exactly), but
  they have no place in a base plane.
\item \textbf{The five polynomial-sector rules} $74$, $88$, $104$, $173$, $233$:
  confirm $b=2$ directly rather than through a fit, e.g.\ by showing
  $D_{\max}/2^N$ is bounded below.
\item \textbf{A basin histogram in the Tier-1a schema}, so the gap of
  \S\ref{sec:basins} cannot recur for future units.
\end{enumerate}

\section{Reproduce}

\begin{verbatim}
# the sector sweep (N in the outer loop: uniform coverage if stopped early)
python -m qca_fragmentation.wcc_sweep --which all --bc obc0 \
       --n-min 6 --n-max 16 --no-stop-on-ergodic --order n

# a single rule
python -m qca_fragmentation.run_rule --rule 156 --N 6-16 --bc obc0 --tier wcc

# recompute the basin units the Tier-1a archive cannot answer
python -m qca_fragmentation.basin_recompute --bc obc0

# descriptors + hyperbola, then tables and figures F1-F5
python -m qca_fragmentation.scaling.sectors      --bc obc0
python -m qca_fragmentation.scaling.r9_tables    --bc obc0
python -m qca_fragmentation.scaling.sector_figure --bc obc0 --rebuild

python -m pytest qca_fragmentation/tests -q
\end{verbatim}

\begin{sloppypar}
Data: \fn{results\_tier1e/\{rule\}\_\{bc\}.jsonl} (one record per unit with the
full sector-size histogram, the derived observables and the A2--A4 check
results), \fn{results\_basins/\{rule\}\_\{bc\}.jsonl} (the independent basin
recomputation), \fn{analytics/sector\_map\_obc0.json} (per-rule descriptors,
hyperbola verdicts and attractor deficits),
\fn{analytics/sector\_validation\_obc0.json}, and
\fn{checkpoints/manifest.jsonl} for the audit trail.
\end{sloppypar}

\end{document}
